\makeglossaries
\newglossaryentry{Hardcoded}{
        name=Hardcoded,
        description={É um termo utilizado para descrever um valor que foi codificado diretamente no código de um programa e que não depende de nenhuma variável ou entrada do usuário}
}

\newglossaryentry{mux}{
        name=MUX,
        description={Abreviatura de Multiplexador. É um circuito digital que seleciona um entre vários sinais de entrada, com base em sinais de controlo, e transmite-o para uma única saída}
}

\newglossaryentry{carry}{
        name=Carry,
        description={É um bit que é transportado de um bit para o seguinte numa operação de soma, indicando que a soma excedeu o limite de representação}
}

\newglossaryentry{cout}{
        name=Carry Out,
        description={É o carry que é transportado para fora do circuito, indicando que a soma excedeu o limite de representação}
}

\newglossaryentry{vhdl}{
        name=VHDL,
        description={VHDL é uma linguagem de descrição de hardware utilizada para modelar, simular e sintetizar circuitos digitais em sistemas eletrónicos. A sigla significa VHSIC Hardware Description Language, onde VHSIC é a abreviatura de Very High Speed Integrated Circuit}
}

\newglossaryentry{clock}{
        name=Clock,
        description={É um sinal periódico que é utilizado para sincronizar as operações de um circuito digital}
}


\newacronym{alu}{ALU}{Arithmetic Logic Unit}
\newacronym{asm}{ASM}{Algorithmic State Machine}
\newacronym{rom}{ROM}{Read Only Memory}
\newacronym{lsd}{LSD}{Lógica de Sistemas Digitais}
\newacronym{ffd}{FFD}{Flip-Flop D}
\newacronym{lsl}{LSL}{Logic Shift Left}
\newacronym{led}{LED}{Light-Emitting Diode}